\section{Pricing retained layers by certificates}\label{sec:certificate-pricing}

Retaining layers in a quotient imposes a lattice condition on the new
column. We now count partial certificates of that condition. The resulting
estimate permits the certificate size to be chosen separately in each
class of outcomes; this will allow a later argument to choose it according
to the largest potential in that class.

We first record the probability bound for a fixed certificate at any
fixed cube precision \(A_{\mathrm{cube}}>0\). Let \(\mathcal F\) be exposed
information that fixes a nonsingular binary matrix \(C\) of size
\(m\in\mathcal I_N\) and a quotient map
\(\pi:\Z^m\twoheadrightarrow G\) with \(C\Z^m\subseteq\ker\pi\).
Suppose that, conditionally on \(\mathcal F\), a column \(v\) is uniform
on \(\bits^m\). Take \(C\) in the typical set of
Proposition~\ref{prop:cube-concentration} at precision
\(A_{\mathrm{cube}}\).

Fix a candidate set \(E\) of positive layers of \(G\) after these data
are fixed, and let \(\mathcal R\) be the layers retained on quotienting
\(G\) by \(\pi(v)\). Lemma~\ref{lem:retention-lattice} gives a lattice
\(K_E\), fixed by \(\pi\) and \(E\), such that \(E\subseteq\mathcal R\)
implies \(v\in K_E\), with
\[
 C\Z^m\subseteq K_E\subseteq\Z^m,\qquad
 [\Z^m:K_E]=N^{Nq(E)},\qquad
 q(E)=\sum_{(p,e)\in E}k_G(p,e)\nu_p.
\]
The group-order bound in Equation~\eqref{eq:group-order} gives
\(0\le q(E)\le1/2\). Conditional uniformity of \(v\) and the
simultaneous cube estimate therefore give
\begin{equation}\label{eq:fixed-certificate-probability}
 \begin{aligned}
 \Pr(E\subseteq\mathcal R\mid\mathcal F)
 &\le \Pr(v\in K_E\mid\mathcal F)
 \le 2^{-m}M_{K_E}\\
 &\le 2^{-m}2^{N(1-c(q(E))+\ell^{-A_{\mathrm{cube}}})}
 =2^{-N(c(q(E))-\eta_N)},\\
 \eta_N&=\ell^{-A_{\mathrm{cube}}}+1-\frac mN.
 \end{aligned}
\end{equation}
Here \(1-m/N\le N^{-3/10}\). The same typical set works for every
exposed quotient, since all its certificate lattices contain the original
\(C\Z^m\).
We will handle marks that depend on \(v\) by counting the fixed candidates
and how many each outcome offers.

For \(0\le t\le1\), let
\[
 H_2(t)=-t\log_2t-(1-t)\log_2(1-t),
 \qquad H_2(0)=H_2(1)=0.
\]
For \(s,u\ge0\) and \(0\le x\le s\), define
\begin{equation}\label{eq:entropy-cost}
 \mathcal E(s,u,x)
 =(s+u)H_2\!\left(\frac{x}{s+u}\right)
       -sH_2\!\left(\frac{x}{s}\right),
\end{equation}
using the continuous values when a denominator vanishes. Here \(s\) and
\(u\) will be the total weights of marked and unmarked layers, and \(x\)
will be a target weight for a certificate chosen among the marked layers.
When all layers have weight \(1/N\), the ratio
\(\binom{(s+u)N}{xN}/\binom{sN}{xN}\) counts possible certificates
relative to those offered by a fixed marked set, and its base-two logarithm
is \(N\mathcal E(s,u,x)+O(\log N)\) for integer counts.

\begin{proposition}[Certificate pricing]\label{prop:certificate-pricing}
Fix \(A>0\) and at most two positive constants \(w_i\). Typically for
nonsingular binary matrices \(C\) of size
\(m\in\mathcal I_N\), the following assertion holds. The typical set
is obtained by imposing Proposition~\ref{prop:cube-concentration} at a
sufficiently large fixed precision depending on \(A\) and these constants.

Let \(\mathcal F\) denote exposed information which fixes \(C\), a
quotient map \(G_C\twoheadrightarrow G\), and the data described next.
Suppose that, conditionally on \(\mathcal F\), a tested column \(v\) is
uniform on \(\bits^m\). Partition the layers of \(G\) into at most two
categories, each a union of whole plateaus. In category \(i\), assume
that every height is at least \(w_i\); only \(w_i=1,2\) will be needed.
The partition is part of the exposed data. In the quotient of \(G\) by the image of \(v\),
an outcome may mark some of the retained layers. Require that the marked
layers on each plateau form a prefix. The marks may also depend on other
random bits in the experiment.

Let \(\mathcal A\) be any subevent of this experiment. There is a
partition of \(\mathcal A\) into at most \(2^{N\ell^{-A}}\) classes,
constructed before choosing a representative or any target weights, with
the following property. Take any nonempty class \(\mathcal A_\alpha\),
any representative outcome in it, and write \(s_i\) and \(u_i\) for
the representative's sums of the weights \(\nu_p\) of the marked and
unmarked layers in category \(i\). For any choices \(0\le x_i\le s_i\),
\begin{equation}\label{eq:certificate-pricing}
 \Pr(\mathcal A_\alpha\mid\mathcal F)
 \le 2^{-N\left(
 c\left(\sum_i w_i x_i\right)
       -\sum_i\mathcal E(s_i,u_i,x_i)-\ell^{-A}\right)}.
\end{equation}
Here \(\mathcal A_\alpha\) already includes the restriction to
\(\mathcal A\); the probability is conditioned only on the exposed
information. The class count and the estimate are uniform in that
information, the quotient, the subevent, and the representative. In
particular, the exposed information may include a preceding column that
fixes \(G\), provided the tested column remains conditionally uniform.
\end{proposition}

We record the entropy facts needed below. Put
\[
 H(s,u)=(s+u)H_2\!\left(\frac{s}{s+u}\right),\qquad H(0,0)=0.
\]
Expansion of the logarithms gives
\begin{equation}\label{eq:entropy-difference}
 \mathcal E(s,u,x)=H(s,u)-H(s-x,u).
\end{equation}
For \(s>0\),
\[
 \frac{\partial H}{\partial s}(s,u)=\log_2(1+u/s),
\]
which is nonnegative and nonincreasing in \(s\). Thus \(\mathcal E\ge0\).
For fixed \(s,u\), the function \(\mathcal E(s,u,x)\) is convex in
\(x\in[0,s]\): in the interior its derivative is
\(\log_2(1+u/(s-x))\), which is nondecreasing in \(x\). For fixed
\(u,x\), it is nonincreasing in \(s\ge x\), since for \(s>x\),
\[
 \frac{\partial\mathcal E}{\partial s}(s,u,x)
 =\log_2(1+u/s)-\log_2(1+u/(s-x))\le0.
\]
The endpoint cases follow by continuity. The defining formula also gives
homogeneity of degree one in all three size variables:
\(\mathcal E(ts,tu,tx)=t\mathcal E(s,u,x)\) for \(t\ge0\).
For fixed \(0\le\theta\le1\), write
\[
 \mathcal E_\theta(s,u)=\mathcal E(s,u,\theta s).
\]
This function is concave and homogeneous of degree one in \((s,u)\).
Indeed, the first term in
\[
 \mathcal E_\theta(s,u)
 =(s+u)H_2\!\left(\frac{\theta s}{s+u}\right)-sH_2(\theta)
\]
is the perspective of the concave function \(H_2\), composed with a
linear map, and the second term is linear. Consequently
\(\mathcal E_\theta\) is superadditive on the nonnegative quadrant:
\begin{equation}\label{eq:entropy-superadditive}
 \sum_j\mathcal E_\theta(s_j,u_j)
 \le \mathcal E_\theta\!\left(\sum_j s_j,\sum_j u_j\right).
\end{equation}

These functions have uniform continuity estimates at their boundaries.
For every fixed \(M\), the function \(-t\log_2t\), extended by zero at
zero, is H\"older continuous of exponent \(1/2\) on \([0,M]\). To see
this, integrate its derivative on an interval of length \(d\); the
absolute integral is at most
\(O_M(d(1+|\log d|))=O_M(\sqrt d)\), with the same conclusion for
intervals meeting zero. Expressing \(H(s,u)\) as
\[
 -s\log_2s-u\log_2u+(s+u)\log_2(s+u)
\]
therefore shows that it is H\"older continuous of exponent \(1/2\) on
every bounded nonnegative rectangle. Equation~\eqref{eq:entropy-difference}
gives the same estimate for \(\mathcal E(s,u,x)\) on bounded parts of
its domain. It also gives an estimate for \(\mathcal E_\theta(s,u)\)
whose constant is uniform over \(0\le\theta\le1\). We will use these
estimates both to round certificate cardinalities and to replace layer
weights by nearby cell weights.

\begin{proof}
Work conditionally on fixed exposed data satisfying the assumptions.
Choose a fixed cube precision \(A_{\mathrm{cube}}>0\), to be specified at
the end, and invoke Equation~\eqref{eq:fixed-certificate-probability} for
each fixed candidate \(E\). Here the quotient map
\(G_C\twoheadrightarrow G\), composed with the natural map from \(\Z^m\),
gives the exposed map \(\pi\) in that observation. This uses only the
conditional uniformity of \(v\); additional random bits and the restriction
to \(\mathcal A\) need not be independent of \(v\).

We now construct the classes. Fix a large constant \(D>0\), to be
chosen at the end, and set \(\delta=\ell^{-D}\) and
\(a_p=\log_Np\). The order bound gives
\begin{equation}\label{eq:certificate-order-bounds}
 \sum_{p,e}k_G(p,e)\frac{a_p}{N}\le\frac12,
 \qquad
 \sum_{p,e}k_G(p,e)\le\frac{N\log N}{2\log2}.
\end{equation}
We classify the marks in three prime ranges.

For \(a_p>\ell^D\), record every mark bit. There are at most
\(N/(2\ell^D)\) layers in this range by the first bound in
Equation~\eqref{eq:certificate-order-bounds}, so this record costs
\(O(N\ell^{-D})\) bits.

For \(a_p<1-\delta\), record the number of marks on each plateau.
These counts determine all the marks because they are prefixes.
There are at most \(N^{1-\delta}\) primes in this range. If \(d_p\)
is the number of plateaus at \(p\), their distinct positive heights
give
\[
 \frac{d_p(d_p+1)}2\le\sum_e k_G(p,e).
\]
Cauchy--Schwarz and the second bound in
Equation~\eqref{eq:certificate-order-bounds} show that the total number
of these plateaus is
\(O(N^{1-\delta/2}\sqrt{\log N})\). Each prefix count costs
\(O(\log N)\) bits, since every plateau length is \(O(N\log N)\).
The total cost is therefore
\[
 O\!\left(N^{1-\delta/2}(\log N)^{3/2}\right)
       =o(N\ell^{-D}).
\]
For the last equality, use
\(\delta\log N=(\log N)/\ell^D\gg\log\log N\).

The remaining primes satisfy \(1-\delta\le a_p\le\ell^D\).
For each category, split this range into bins of width at most
\(\delta\), starting at \(1-\delta\), and call a category-bin pair
a cell. There are \(J=O(\ell^{2D})\) cells. In a cell \(j\), let
\(\lambda_j\) be the lower endpoint of its bin, so
\[
 1-\delta\le\lambda_j\le\ell^D,
 \qquad \lambda_j\le a_p\le\lambda_j+\delta
 \quad\hbox{for its layers}.
\]
The number of layers in this entire range is \(O(N)\), again by the
first order bound. Record only the number of marks in each cell. This
costs \(O(\ell^{2D}\log N)=o(N\ell^{-D})\) bits. All these records
depend only on the marks and the exposed data; intersecting their classes
with \(\mathcal A\) gives the desired classes for the subevent. Their
number is at most \(2^{O(N\ell^{-D})+o(N\ell^{-D})}\), uniformly in
the exposed quotient.

Fix one class, a representative, and target weights as in the statement.
For category \(i\), put \(\theta_i=x_i/s_i\) if \(s_i>0\), and put
\(\theta_i=0\) otherwise. In the large and small prime ranges the
marked sets are known throughout the class; include all of them in each
certificate. In an intermediate cell \(j\) of category \(i\), let
\(n_j\) be its number of layers and \(r_j\) its recorded number of
marks, and require a certificate to contain any
\[
 a_j=\lfloor\theta_i r_j\rfloor
\]
of those marked layers. Thus every outcome in the class offers exactly
\[
 \mathfrak m=\prod_j\binom{r_j}{a_j}
\]
certificates, whereas the number of possible certificates for the class
is at most
\[
 \mathfrak n=\prod_j\binom{n_j}{a_j}.
\]
Empty products have value one. The latter count is an upper bound because
it allows arbitrary layers in each cell, without imposing the prefix
restriction.

We check the weight certified by each such choice uniformly over all
outcomes in the class. In category \(i\), let \(s_i^0\) be the marked
weight in the large and small prime ranges; it is fixed in the class.
The representative's marked weight in its intermediate cells differs
from \(\sum_{j\in i}\lambda_jr_j/N\) by at most
\(\delta\sum_{j\in i}r_j/N=O(\delta)\). The weight selected in those
cells is at least
\[
 \sum_{j\in i}\frac{\lambda_j a_j}{N}
 \ge \theta_i\sum_{j\in i}\frac{\lambda_j r_j}{N}
              -\frac1N\sum_{j\in i}\lambda_j.
\]
The last term is \(O(\ell^{3D}/N)=o(\delta)\). Since all of the
known marked weight is included and \(s_i^0\ge\theta_i s_i^0\), the
selected weight in this category is at least \(x_i-O(\delta)\).
Every selected layer there has height at least \(w_i\). Consequently,
uniformly over its choices,
\begin{equation}\label{eq:certificate-index-lower-bound}
 q(E)\ge q_0-O(\delta),\qquad q_0=\sum_i w_i x_i.
\end{equation}
Both \(q(E)\) and \(q_0\) belong to \([0,1/2]\): for the latter,
\(q_0\le\sum_i w_i s_i\le\sum_{p,e}k_G(p,e)\nu_p\).
The profile \(c\) is increasing and H\"older continuous of exponent
\(1/2\) on this interval. Equation~\eqref{eq:certificate-index-lower-bound}
therefore gives
\[
 c(q(E))\ge c(q_0)-O(\sqrt\delta).
\]
Together with Equation~\eqref{eq:fixed-certificate-probability}, this
bounds the probability of offering any one fixed candidate certificate
by \(2^{-N(c(q_0)-\eta_N-O(\sqrt\delta))}\). This conclusion holds
after intersection with the class and with \(\mathcal A\), since such
an intersection still implies \(v\in K_E\).

Double counting outcome-certificate pairs now yields
\begin{equation}\label{eq:certificate-double-count}
 \Pr(\mathcal A_\alpha\mid\mathcal F)
 \le \frac{\mathfrak n}{\mathfrak m}
       2^{-N(c(q_0)-\eta_N-O(\sqrt\delta))}.
\end{equation}
Indeed, summing over candidate certificates the probabilities of the
outcomes in the class offering that certificate gives
\(\mathfrak m\Pr(\mathcal A_\alpha\mid\mathcal F)\); each summand
is bounded by the fixed-certificate estimate. This also explains why
no conditional independence of the marks or of the subevent is needed.

It remains to bound the ratio in Equation~\eqref{eq:certificate-double-count}.
For integers \(0\le a\le n\), the elementary binomial bounds are
\[
 nH_2(a/n)-\log_2(n+1)
 \le\log_2\binom na\le nH_2(a/n).
\]
For completeness, insert the parameter \(a/n\) into the binomial
distribution. The upper bound follows because its probability at \(a\)
is at most one, and the lower bound follows because \(a\) is a mode and
there are \(n+1\) probabilities summing to one. The endpoints are
immediate. Applying these bounds to a cell gives
\[
 \frac1N\log_2\frac{\binom{n_j}{a_j}}{\binom{r_j}{a_j}}
 \le
 \mathcal E\!\left(\frac{r_j}{N},\frac{n_j-r_j}{N},\frac{a_j}{N}\right)
       +\frac{\log_2(r_j+1)}N.
\]
All arguments here stay in a fixed bounded set. Replacing \(a_j/N\)
by \(\theta_i r_j/N\) changes the entropy term by \(O(N^{-1/2})\),
uniformly in \(\theta_i\), by the H\"older estimate above. After summing
over \(J=O(\ell^{2D})\) cells, all rounding and binomial errors are
\[
 O\!\left(\ell^{2D}\left(N^{-1/2}+\frac{\log N}{N}\right)\right),
\]
which is smaller than every fixed negative power of \(\ell\).

We convert the remaining entropy terms from counts to weights. Their
nonnegativity and homogeneity give, for a cell in category \(i\),
\[
 \mathcal E_{\theta_i}\!\left(\frac{r_j}{N},\frac{n_j-r_j}{N}\right)
 \le\frac1{1-\delta}
 \mathcal E_{\theta_i}\!\left(
       \frac{\lambda_jr_j}{N},\frac{\lambda_j(n_j-r_j)}{N}\right).
\]
Let \(u_i^0\), alongside \(s_i^0\), denote the unmarked weight in the
known large and small ranges, and put
\[
 \underline s_i=s_i^0+\sum_{j\in i}\frac{\lambda_jr_j}{N},
 \qquad
 \underline u_i=u_i^0+\sum_{j\in i}\frac{\lambda_j(n_j-r_j)}{N}.
\]
By Equation~\eqref{eq:entropy-superadditive} the weighted cell terms
sum to at most
\(\mathcal E_{\theta_i}(\underline s_i,\underline u_i)\).
Here adding the known ranges preserves this upper bound because their
entropy term is nonnegative. The representative weights satisfy
\[
 0\le s_i-\underline s_i=O(\delta),\qquad
 0\le u_i-\underline u_i=O(\delta).
\]
All these weights have bounded sum, at most \(1/2\) before replacing
weights by lower endpoints. The uniform H\"older estimate for
\(\mathcal E_{\theta_i}\) and the factor
\((1-\delta)^{-1}=1+O(\delta)\) therefore show that
\begin{equation}\label{eq:certificate-entropy-bound}
 \frac1N\log_2\frac{\mathfrak n}{\mathfrak m}
 \le \sum_i\mathcal E(s_i,u_i,x_i)+O(\sqrt\delta)
   +O\!\left(\ell^{2D}\left(N^{-1/2}+\frac{\log N}{N}\right)\right).
\end{equation}
The constants in this estimate are uniform in the targets, including
their endpoint values.

Choose the fixed \(D\) sufficiently large that
\(\ell^{-D/2}=o(\ell^{-A})\), and then choose a fixed cube precision
\(A_{\mathrm{cube}}>A\). The class count is at most \(2^{N\ell^{-A}}\) for large
\(N\). Moreover, \(\eta_N\), the rounding errors, and the two
\(O(\sqrt\delta)\) errors in
Equations~\eqref{eq:certificate-double-count} and
\eqref{eq:certificate-entropy-bound} have sum at most \(\ell^{-A}\)
for large \(N\). Substitution proves
Equation~\eqref{eq:certificate-pricing}. Every choice of precision was
fixed before \(N\) tended to infinity and was independent of the
quotient and exposed columns, as required.
\end{proof}

We next extract a uniform cost for retaining a noticeable total weight.
Set
\begin{equation}\label{eq:retention-scales}
 R=\lceil\ell^2\rceil,\qquad \sigma=R^{-6}.
\end{equation}

\begin{proposition}[Retained-layer cost]\label{prop:retention-cost}
Typically for nonsingular binary matrices \(C\) of size
\(m\in\mathcal I_N\), the following estimate holds. Let \(h\in\{1,2\}\)
and let \(v_1,\ldots,v_h\) be independent uniform columns in
\(\bits^m\), conditionally on \(C\). Starting with \(G^{(0)}=G_C\),
let \(G^{(j)}\) be the quotient of \(G^{(j-1)}\) by the image of
\(v_j\). For each \(1\le j\le h\), and conditionally on any values
of the preceding columns,
\begin{equation}\label{eq:retention-cost}
 \Pr\!\left(
  \sum_{\substack{(p,e)\text{ retained in}\\
                  G^{(j-1)}\twoheadrightarrow G^{(j)}}}\nu_p
       \ge\sigma
       \,\middle|\, C,v_1,\ldots,v_{j-1}\right)
 \le 2^{-N\sigma/2}.
\end{equation}
The typical set imposes one fixed precision on \(C\), independent of
the values of the preceding columns. When these columns occur in a
random enlargement of \(C\), no nonsingularity assumption on the
destination matrix is required.
\end{proposition}

\begin{proof}
Apply Proposition~\ref{prop:certificate-pricing} at the fixed precision
\(A_0=40\). For the \(j\)-th quotient, condition on \(C\) and the
preceding columns; these fix the quotient map to \(G^{(j-1)}\), while
\(v_j\) remains uniform. Use one category of all its layers, take
\(w_1=1\), and mark all retained layers. Lemma~\ref{lem:layer-quotient}
ensures that these marks are prefixes on plateaus. Restrict to the
subevent in Equation~\eqref{eq:retention-cost}. For any representative
of a nonempty class its marked weight \(s\) is at least \(\sigma\),
so we may choose \(x=\sigma^2\). The total layer weight \(s+u\) is
at most \(1/2\), by Equation~\eqref{eq:group-order}. Since
\(H_2(t)\le t\log_2(e/t)\), Equation~\eqref{eq:entropy-cost} yields
\[
 0\le\mathcal E(s,u,\sigma^2)
 \le (s+u)H_2\!\left(\frac{\sigma^2}{s+u}\right)
 \le O\!\left(\sigma^2\log(1/\sigma)\right).
\]
This bound is uniform over the representatives in this subevent.

The defining formula for the profile \(c\) satisfies
\(c(q)=2\sqrt q+O(q)\) as \(q\downarrow0\). Indeed, in that formula
\(u=2q/(1-2q)=2q+O(q^2)\), so
\(\sqrt{2u-u^2}=2\sqrt q+O(q^{3/2})\), and substitution gives the
claimed expansion. In particular \(c(q)\ge(3/2)\sqrt q\) for all
sufficiently small \(q\). Proposition~\ref{prop:certificate-pricing}
therefore bounds the probability of each class by
\[
 2^{-N\left((3/2)\sigma
       -O(\sigma^2\log(1/\sigma))-\ell^{-40}\right)}.
\]
There are at most \(2^{N\ell^{-40}}\) classes. Since
\(\sigma\asymp\ell^{-12}\), both
\(\sigma^2\log(1/\sigma)=o(\sigma)\) and
\(\ell^{-40}=o(\sigma)\). Summing the class probabilities gives
Equation~\eqref{eq:retention-cost} for large \(N\).

The argument tests only the quotient by the selected column. It does not
use a new principal block or any property of an enlarged matrix. The
conditional uniformity for a second column and the simultaneous cube
estimate for lattices containing \(C\Z^m\) establish the claimed
uniformity over the preceding column.
\end{proof}
